% !TeX root = ../PIFB2.tex
\documentclass[11pt]{article}
\usepackage[utf8]{inputenc}
\usepackage{amsmath,amssymb,amsfonts,amsthm}
\usepackage{graphicx}
\usepackage{xcolor}
\usepackage{longtable}
\usepackage{hyperref}
\usepackage{tikz}
\usetikzlibrary{calc, arrows.meta, decorations.markings, patterns, shapes.geometric, positioning, 3d, decorations.pathmorphing}
\usepackage{geometry}
\usepackage{booktabs}
\PassOptionsToPackage{hyphens}{url}
\usepackage{tabularx}
\geometry{margin=1in}
\emergencystretch=3em
\Urlmuskip=0mu plus 1mu\relax
\def\UrlBreaks{\do\/\do\-\do\.\do\_}
\newenvironment{acknowledgments}
  {\section*{Acknowledgments}}
  {}
\usepackage{cite}
\usepackage{microtype}
\hypersetup{hidelinks,pdftitle={Participatory It from Bit},pdfauthor={Robert C. Dennis}}

% Theorem environments
\newtheorem{theorem}{Theorem}
\newtheorem{lemma}{Lemma}
\newtheorem{proposition}{Proposition}
\newtheorem{corollary}{Corollary}
\newtheorem{definition}{Definition}
\newtheorem{remark}{Remark}
\newtheorem{postulate}{Postulate}
\newtheorem{assumption}{Assumption}

\providecommand{\KL}{\operatorname{KL}}
\providecommand{\softmax}{\operatorname{softmax}}

\begin{document}

\title{Participatory It from Bit: Selective Communication, Inference, and Effective Physical Structure}

\author{
  Robert C. Dennis\thanks{Corresponding author. Email: cdenn016@gmail.com}\\[0.3cm]
  \small Independent Researcher\\
  \small Leander, Texas, United States\\[0.2cm]
  \small ORCID: \href{https://orcid.org/0009-0004-4688-2341}{0009-0004-4688-2341}
}

\date{\today}

\maketitle

\begin{abstract}
Wheeler's participatory it-from-bit program motivates a theory of how local records, communication and learned descriptions can support a shared physical world. This monograph develops a classical probabilistic framework in which agents are local sections of an associated statistical bundle over a continuous, possibly curved base. At every point of overlap, a network of selective matrix-valued channels communicates only part of each agent's information. A normalized population model yields an exact variational free energy; a separate channel-consensus objective describes agreement when no such generative identification has been established. One joint statistical fiber accommodates both current-state beliefs and uncertain models. Independent invertible interactions are represented by latent Lie-algebra generators with a normalized Gaussian-based loop prior. Their recognition flow is first-order Fisher descent; a separately specified episode-history prior gives Euclidean history inertia and a controlled Proca-type linearized link equation. Prediction and conditional expectation derive cross-scale priors and messages once a process law and retained variables are specified. Exact coarse-graining produces effective interactions and often memory, while a conditional four-agent selection theorem illustrates what formation dynamics can establish. Worked constructions distinguish Fisher stiffness, dissipative response and metastable persistence from mechanical mass. Information clocks and collective distinguishability are developed with explicit covariance and gauge conventions. A complex two-dimensional message sector supplies the full proper orthochronous Lorentz group through Hermitian determinant geometry; identifying this internal cone with physical spacetime requires a further soldering structure and dynamics. The resulting program minimizes independent constructions while retaining the assumptions that normalization, inference and coarse-graining cannot remove. It establishes conditional mathematical mechanisms, not a derivation of quantum theory, gravity or matter from information alone.
\end{abstract}

\noindent\textbf{Keywords:} Participatory realism $\cdot$ It from bit $\cdot$ Gauge theory $\cdot$ Information geometry $\cdot$ Variational free energy $\cdot$ Multi-agent inference

\clearpage
\setcounter{tocdepth}{1}
\tableofcontents
\clearpage

\part*{Part I: The Construction}\addcontentsline{toc}{part}{Part I: The Construction}

